\begin{theorem}[Wild quadratic boundary cancellation]\label{mod:cases-wildquadratic-boundary}
Let the exact wild-quadratic stationary assembly lie in its exceptional
boundary row.  Lean first recovers the boundary condition from the retained
coefficient view itself:
\[
  m=t+1,\qquad b=m-1=t,\qquad T=t+1\text{ is odd}.
\]
Thus the condition is not replaced by an unrelated parity hypothesis.

The input is the exact \texttt{WildQuadraticErrorFormula}, together with the
two actual refinement-correction records selected by that formula.  These
records retain their admissible denominators, stationary numerators,
normalized quotient classes, affine coefficients, and positively oriented
inverse-character identities.  The representative translations used
upstream have the precise direction
\[
 H_{\mathrm{after}}=H_{\mathrm{before}}(\,\cdot+a),\qquad
 E_{\mathrm{after}}=H_{\mathrm{before}}(a)E_{\mathrm{before}},
\]
so the elementary value and residual Hasse phase are transported together.
Consequently the complete quotient, in the numerator/denominator orientation
\[
 \frac{L_{\chi_K}L_\tau}{L_{\chi_F}L_{\tau\chi_F}},
\]
is independent of the chosen compatible representatives.

Put
\[
 A_\rho=\frac{\rho}{1+\rho},\qquad
 B=\frac{\rho^2}{1+\rho^2},\qquad
 C=\frac{\rho}{1+\rho^2},\qquad
 \gamma_1=\frac{\gamma_0+\rho\gamma+r}{1+\rho},
 \quad r^2=\rho.
\]
The exact error formula, including both positive quadratic correction factors
and the translated elementary phase, is
\[
 \operatorname{Err}_{K/F}(\chi_F)
 =\mathfrak q(A_\rho)\mathfrak q(B)\mathfrak q(C)\psi_0(a_1),
 \qquad
 a_1=\gamma_1A_\rho+\gamma B+\gamma_0C+C+A_\rho.
\]
No affine coefficient, polar sign, correction unit, or representative
translation is discarded in passing to this formula.

The proof splits on the characteristic of $F$.  In equal characteristic
two, the exact lower-break theorem gives $t$ odd, while the boundary row
gives $t+1$ odd; hence this row is impossible.  In mixed characteristic,
the positive-polar refinement identities give
\[
 B+C=A_\rho,
 \qquad
 \mathfrak q(B)\mathfrak q(C)
   =\mathfrak q(A_\rho)\psi_0(BC),
 \qquad
 \mathfrak q(A_\rho)^2=\psi_0(A_\rho).
\]
All remaining terms therefore form the exact residue exponent
\[
 a_1+A_\rho+BC.
\]
With the manuscript's numerator and denominator orientation, define
\[
 z=\frac{r^3}{1+\rho^2}+\frac{1}{1+\rho}.
\]
A direct characteristic-two calculation, using the nonvanishing of
$1+\rho$ and $1+\rho^2$, proves
\[
 a_1+A_\rho+BC=z+z^2.
\]
The Artin--Schreier sign is the plus sign shown here.  Hence
\(\psi_0(z+z^2)=1\), and the complete boundary error is exactly
\[
 \operatorname{Err}_{K/F}(\chi_F)=1.
\]

\LeanFile{LanglandsFirstMainLemma/Cases/WildQuadratic/Boundary.lean}
\lean{LanglandsFirstMainLemma.WildQuadraticBoundaryData.artinSchreierWitness}
\lean{LanglandsFirstMainLemma.WildQuadraticBoundaryData.B_add_C_eq_A_rho}
\lean{LanglandsFirstMainLemma.WildQuadraticBoundaryData.boundary_residual_eq_artinSchreier}
\lean{LanglandsFirstMainLemma.WildQuadraticBoundaryData.value_eq_one}
\lean{LanglandsFirstMainLemma.WildQuadraticBoundaryFormula.exceptionalCondition}
\lean{LanglandsFirstMainLemma.WildQuadraticBoundaryFormula.equalCharacteristic_elim}
\lean{LanglandsFirstMainLemma.wildQuadratic_boundary}
\leanok
\SourceText{Proposition prop:wild-quadratic-cancellation, boundary case.}
\uses{mod:cases-wildquadratic-errorformula,mod:cases-wildquadratic-quadraticrefinement,mod:cases-wildquadratic-equalcharacteristicbreak,mod:parameters-minimalorbitstationary}
\FormalizationContract
\end{theorem}
